% META: {"id":"1SPE-DERIVATION-LOCAL-CR-011","chapitre":"1SPE-DERIVATION-LOCAL","type_objet":"cours","capacites_codes":["C2"],"sources_inspiration":[],"status":"generated"}
\section{Du taux de variation au nombre dérivé}
\definition[Nombre dérivé]{Soit $f$ une fonction et $a$ un réel. Lorsque les taux de variation de $f$ entre $a$ et $a+h$ se rapprochent d'un même nombre lorsque $h$ est très proche de $0$, ce nombre est le \textbf{nombre dérivé} de $f$ en $a$. On le note $f'(a)$.}
\exemple{Pour la fonction carré, au voisinage de $2$, les taux de variation $\frac{f(2+h)-f(2)}h$ sont proches de $4$ lorsque $h$ est proche de $0$. On retient $f'(2)=4$.}
\contreexemple{Le nombre $f'(2)$ ne se calcule pas en remplaçant $h$ par $0$ dans le quotient : ce quotient n'est alors pas défini.}
\propriete[Interprétation géométrique]{Si $f'(a)$ existe, il est le coefficient directeur de la tangente à la courbe de $f$ au point d'abscisse $a$.}
\exemple{Si $f'(3)=-2$, la tangente descend de $2$ unités quand on avance d'une unité : sa pente est négative.}
\contreexemple{Une pente négative ne signifie pas que toutes les images de $f$ sont négatives ; elle indique seulement un comportement local de la courbe.}
\begin{nxfigure}{Sécantes convergeant vers la tangente}
\begin{tikzpicture}
\begin{axis}[nexus,
  xmin=-0.5, xmax=5, ymin=-0.5, ymax=10,
  xlabel=$x$, ylabel=$y$,
  xtick={0,1,2,3,4}, ytick={0,2,4,6,8}]
  \addplot[courbe,domain=0:4.5,samples=60] {0.5*x^2};
  % Sécantes depuis a=2
  \addplot[secante,domain=0.5:4.5] {0.5*(x-2)*(2+4) + 2};  % h=2, pente=3
  \addplot[secante,domain=0.5:4.5] {2.5*(x-2) + 2};        % h=1, pente=2.5
  \addplot[secante,domain=0.5:4.5,nxOr!80] {2.25*(x-2) + 2}; % h=0.5
  % Tangente
  \addplot[tangente,domain=0.5:4] {2*(x-2) + 2};
  % Point de contact
  \addplot[point] coordinates {(2,2)};
  \node[anchor=north west,font=\footnotesize,nxBleu] at (axis cs:2,2) {$A(a;f(a))$};
  \node[anchor=south,font=\footnotesize,nxOr] at (axis cs:3.8,5.6) {\small tangente};
\end{axis}
\end{tikzpicture}
\nxcaption{Les sécantes (pointillés) convergent vers la tangente (trait plein or) lorsque $h\to 0$.}
\end{nxfigure}

\margeAppui{On passe des sécantes à la tangente : les sécantes utilisent deux points, la tangente décrit la direction de la courbe en un seul point.}
\margeAppui{En cinématique, une vitesse moyenne se calcule sur un intervalle ; une vitesse instantanée est interprétée à un instant précis.}
\erreurFrequente{\textbf{Copie fautive :} « $f'(a)=f(a)$. » \textbf{Correction :} $f(a)$ est l'ordonnée du point de la courbe ; $f'(a)$ est la pente de la tangente en ce point.}
\erreurFrequente{\textbf{Copie fautive :} « La vitesse instantanée est la distance parcourue. » \textbf{Correction :} si $d$ est une distance, $d'(t)$ est une vitesse ; l'unité est par exemple m/s.}
\approfondissement{Pour $f(x)=x^2$, $\frac{f(a+h)-f(a)}h=\frac{(a+h)^2-a^2}h=2a+h$ pour $h\ne0$. Ces nombres sont aussi proches que l'on veut de $2a$ lorsque $h$ se rapproche de $0$. Ainsi $f'(a)=2a$ ; ce calcul éclaire l'intuition sans imposer une définition formelle de limite.}
